\documentclass[10pt,aspectratio=169]{beamer}

\usetheme{Madrid}
\usecolortheme{default}
\usepackage{amsmath,amssymb,booktabs,array,multirow}

\definecolor{hkustblue}{RGB}{0,56,116}
\definecolor{hkustgold}{RGB}{166,124,0}
\setbeamercolor{structure}{fg=hkustblue}
\setbeamercolor{frametitle}{fg=white,bg=hkustblue}
\setbeamercolor{title}{fg=white,bg=hkustblue}
\setbeamercolor{block title}{fg=white,bg=hkustblue}
\setbeamercolor{block body}{bg=blue!3}
\setbeamertemplate{navigation symbols}{}
\setbeamertemplate{footline}[frame number]

\newcommand{\pay}[2]{#1,\ #2}

\title[Game Theory Tutorial 1/2]{Tutorial: Introduction to Game Theory (1/2)}
\subtitle{IEDA 1250 Optimizing Decisions for Personal and Business Development}
\author{Wu Hongfan}
\institute{The Hong Kong University of Science and Technology}
\date{June 26, 2026}

\begin{document}

\begin{frame}
  \titlepage
\end{frame}

\begin{frame}{Today's Scope}
\begin{block}{Tutorial 1/2}
Today is the first game theory tutorial. We focus on intuition, payoff tables, Nash equilibrium, dominated strategies, and pure-strategy zero-sum games.
\end{block}
\vspace{0.2cm}
\begin{itemize}
  \item What is a game?
  \item How do we read a payoff matrix?
  \item What does Nash equilibrium mean?
  \item How do we simplify games using dominated strategies?
  \item How do we check for a saddle point in zero-sum games?
\end{itemize}
\vspace{0.2cm}
\begin{alertblock}{Next tutorial}
Mixed strategies, graphical method, LP formulation, and Cournot competition can be continued in Game Theory (2/2).
\end{alertblock}
\end{frame}

\begin{frame}{What Makes a Situation a Game?}
\begin{block}{A game}
A decision situation where each player's payoff depends on the actions of more than one decision maker.
\end{block}
\vspace{0.2cm}
\begin{columns}[T]
\column{0.48\textwidth}
\textbf{Ingredients}
\begin{itemize}
  \item Players
  \item Strategies
  \item Payoffs
  \item Information and timing
\end{itemize}
\column{0.48\textwidth}
\textbf{Typical assumptions}
\begin{itemize}
  \item Rationality
  \item Self-interest
  \item Common knowledge of the game
  \item Strategic interaction
\end{itemize}
\end{columns}
\end{frame}

\begin{frame}{Why Game Theory?}
\begin{itemize}
  \item In LP and IP, one decision maker optimizes under constraints.
  \item In game theory, several decision makers optimize at the same time.
  \item My best decision depends on what you do.
  \item Your best decision depends on what I do.
\end{itemize}
\vspace{0.3cm}
\begin{block}{Examples}
Pricing competition, auctions, political campaigns, traffic routing, and location choices by competing firms.
\end{block}
\end{frame}

\begin{frame}{Example: Prisoner's Dilemma}
\begin{center}
\begin{tabular}{c c c c}
 & & \multicolumn{2}{c}{Prisoner B}\\
 & & Silent & Betray\\
\cmidrule{3-4}
\multirow{2}{*}{Prisoner A} & Silent & $\pay{-2}{-2}$ & $\pay{-10}{0}$\\
 & Betray & $\pay{0}{-10}$ & $\pay{-5}{-5}$\\
\end{tabular}
\end{center}
\vspace{0.3cm}
\begin{itemize}
  \item First number: payoff to Prisoner A.
  \item Second number: payoff to Prisoner B.
  \item Larger payoff is better, so $0>-2>-5>-10$.
\end{itemize}
\end{frame}

\begin{frame}{Best Response Thinking}
\begin{block}{Best response}
A strategy is a best response if it gives the highest payoff given the other player's strategy.
\end{block}
\vspace{0.2cm}
\begin{itemize}
  \item If B is silent, A compares $-2$ with $0$, so A prefers Betray.
  \item If B betrays, A compares $-10$ with $-5$, so A prefers Betray.
  \item By symmetry, B also prefers Betray in both cases.
\end{itemize}
\vspace{0.2cm}
\begin{block}{Conclusion}
Both players choose Betray.
\end{block}
\end{frame}

\begin{frame}{Nash Equilibrium}
\begin{block}{Nash equilibrium}
A strategy profile where no single player can improve by changing only their own strategy.
\end{block}
\vspace{0.2cm}
\begin{itemize}
  \item ``I am doing my best given what you do.''
  \item ``You are doing your best given what I do.''
  \item It is a stability concept, not necessarily a fairness or efficiency concept.
\end{itemize}
\vspace{0.25cm}
\begin{block}{Prisoner's Dilemma}
The Nash equilibrium is $(\text{Betray}, \text{Betray})$.
\end{block}
\end{frame}

\begin{frame}{Equilibrium Can Be Inefficient}
\begin{center}
\begin{tabular}{c c c c}
 & & \multicolumn{2}{c}{Prisoner B}\\
 & & Silent & Betray\\
\cmidrule{3-4}
\multirow{2}{*}{Prisoner A} & Silent & $\pay{-2}{-2}$ & $\pay{-10}{0}$\\
 & Betray & $\pay{0}{-10}$ & $\pay{-5}{-5}$\\
\end{tabular}
\end{center}
\vspace{0.25cm}
\begin{itemize}
  \item Both staying silent is better for both prisoners than both betraying.
  \item But if one prisoner is silent, the other wants to betray.
  \item So $(\text{Silent},\text{Silent})$ is good, but not stable.
\end{itemize}
\end{frame}

\begin{frame}{Example: Beach Kiosk Game}
\begin{itemize}
  \item A beach is 200 meters long.
  \item Two vendors choose locations.
  \item Customers go to the nearest vendor.
  \item Each vendor wants to attract as many customers as possible.
\end{itemize}
\vspace{0.3cm}
\begin{block}{Question}
If the vendors choose locations strategically, where will they locate?
\end{block}
\end{frame}

\begin{frame}{Beach Kiosk Game: Equilibrium Intuition}
\begin{itemize}
  \item If the vendors split the beach evenly at 50m and 150m, each serves half the market.
  \item But each vendor has an incentive to move closer to the center.
  \item The stable outcome is that both vendors locate at the center.
\end{itemize}
\vspace{0.3cm}
\begin{block}{Lesson}
A Nash equilibrium can be stable for the firms but inconvenient for customers. Stability does not imply social efficiency.
\end{block}
\end{frame}

\begin{frame}{How to Read a Payoff Matrix}
\begin{columns}[T]
\column{0.52\textwidth}
\begin{center}
\begin{tabular}{c c c c}
 & & \multicolumn{2}{c}{Column player}\\
 & & L & R\\
\cmidrule{3-4}
\multirow{2}{*}{Row player} & U & $\pay{3}{2}$ & $\pay{0}{3}$\\
 & D & $\pay{2}{1}$ & $\pay{1}{0}$\\
\end{tabular}
\end{center}
\column{0.45\textwidth}
\begin{itemize}
  \item Row player chooses a row.
  \item Column player chooses a column.
  \item First number: row player's payoff.
  \item Second number: column player's payoff.
\end{itemize}
\end{columns}
\vspace{0.25cm}
\begin{alertblock}{Common mistake}
Do not compare the two numbers in the same cell. Compare one player's own payoffs while holding the other player's action fixed.
\end{alertblock}
\end{frame}

\begin{frame}{Finding Nash Equilibrium in a 2-by-2 Game}
\begin{center}
\begin{tabular}{c c c c}
 & & \multicolumn{2}{c}{Player 2}\\
 & & L & R\\
\cmidrule{3-4}
\multirow{2}{*}{Player 1} & U & $\pay{3}{2}$ & $\pay{0}{0}$\\
 & D & $\pay{1}{1}$ & $\pay{2}{3}$\\
\end{tabular}
\end{center}
\vspace{0.2cm}
\begin{enumerate}
  \item For each column, mark Player 1's best row.
  \item For each row, mark Player 2's best column.
  \item A cell with both players' best responses is a Nash equilibrium.
\end{enumerate}
\end{frame}

\begin{frame}{2-by-2 Game: Solution}
\begin{center}
\begin{tabular}{c c c c}
 & & \multicolumn{2}{c}{Player 2}\\
 & & L & R\\
\cmidrule{3-4}
\multirow{2}{*}{Player 1} & U & $\pay{\mathbf{3}}{\mathbf{2}}$ & $\pay{0}{0}$\\
 & D & $\pay{1}{1}$ & $\pay{\mathbf{2}}{\mathbf{3}}$\\
\end{tabular}
\end{center}
\vspace{0.2cm}
\begin{itemize}
  \item If Player 2 chooses L, Player 1 prefers U.
  \item If Player 2 chooses R, Player 1 prefers D.
  \item If Player 1 chooses U, Player 2 prefers L.
  \item If Player 1 chooses D, Player 2 prefers R.
\end{itemize}
\begin{block}{Result}
There are two pure-strategy Nash equilibria: $(U,L)$ and $(D,R)$.
\end{block}
\end{frame}

\begin{frame}{Dominated Strategies}
\begin{block}{Dominated strategy}
Strategy $s$ is dominated if another strategy is always at least as good and sometimes strictly better, regardless of the opponent's action.
\end{block}
\vspace{0.2cm}
\begin{itemize}
  \item For a row player who maximizes: larger row payoffs are better.
  \item For a column player in a zero-sum payoff table for Player 1: smaller entries are better.
  \item Eliminating dominated strategies can simplify a game before solving it.
\end{itemize}
\end{frame}

\begin{frame}{Exercise: Eliminate Dominated Strategies}
\begin{center}
\begin{tabular}{c c c c c c}
 & \multicolumn{5}{c}{Player 2}\\
Player 1 & S1 & S2 & S3 & S4 & S5\\
\midrule
S1 & 3 & 2 & 0 & 3 & 0\\
S2 & 5 & 0 & 2 & 5 & 9\\
S3 & 7 & 3 & 9 & 5 & 9\\
S4 & 4 & 6 & 8 & 7 & 5.5\\
S5 & 6 & 0 & 8 & 8 & 3\\
\end{tabular}
\end{center}
\vspace{0.2cm}
\begin{block}{Task}
Assume the entries are payoffs to Player 1. Player 1 maximizes, and Player 2 minimizes. Which rows or columns can we remove?
\end{block}
\end{frame}

\begin{frame}{Exercise: One Elimination Order}
\begin{enumerate}
  \item Column S4 is dominated by Column S2.
  \item Row S1 is dominated by Row S3.
  \item Row S2 is dominated by Row S3.
  \item In the reduced table, Row S5 is dominated by Row S3.
  \item In the reduced table, Column S3 is dominated by Column S5.
\end{enumerate}
\vspace{0.2cm}
\begin{block}{Takeaway}
Dominance is a simplification method. The order may differ, but each removal must be justified using the remaining table.
\end{block}
\end{frame}

\begin{frame}{Zero-Sum Games}
\begin{block}{Zero-sum game}
One player's gain is exactly the other player's loss.
\end{block}
\vspace{0.2cm}
\begin{itemize}
  \item A single payoff table is enough: entries are payoffs to Player A.
  \item Player A wants to maximize the payoff.
  \item Player B wants to minimize Player A's payoff.
  \item The value of the game is the expected payoff to Player A under optimal play.
\end{itemize}
\end{frame}

\begin{frame}{Saddle Point Test}
\begin{center}
\begin{tabular}{c c c c c c}
 & B1 & B2 & B3 & B4 & Row min\\
\midrule
A1 & 20 & 15 & 12 & 35 & 12\\
A2 & 25 & 14 & 8 & 10 & 8\\
A3 & 40 & 2 & 10 & 5 & 2\\
A4 & -5 & 4 & 11 & 0 & -5\\
\midrule
Col max & 40 & 15 & 12 & 35 & \\
\end{tabular}
\end{center}
\vspace{0.2cm}
\begin{itemize}
  \item Player A chooses the row with largest row minimum: $\max \min = 12$.
  \item Player B chooses the column with smallest column maximum: $\min \max = 12$.
  \item Since they are equal, there is a saddle point.
\end{itemize}
\end{frame}

\begin{frame}{Saddle Point Result}
\begin{block}{Result}
The saddle point is $(A1,B3)$, and the value of the game is $12$.
\end{block}
\vspace{0.2cm}
\begin{itemize}
  \item It is the minimum in Row A1.
  \item It is the maximum in Column B3.
  \item If A announces A1, B cannot reduce A below 12.
  \item If B announces B3, A cannot increase above 12.
\end{itemize}
\end{frame}

\begin{frame}{When There Is No Saddle Point}
\begin{center}
\begin{tabular}{c c c c c}
 & B1 & B2 & B3 & Row min\\
\midrule
A1 & 0 & -2 & 2 & -2\\
A2 & 5 & 4 & -3 & -3\\
A3 & 2 & 3 & -4 & -4\\
\midrule
Col max & 5 & 4 & 2 & \\
\end{tabular}
\end{center}
\vspace{0.2cm}
\begin{itemize}
  \item $\max \min = -2$.
  \item $\min \max = 2$.
  \item They are not equal, so there is no pure-strategy saddle point.
  \item This motivates mixed strategies, which can be covered next time.
\end{itemize}
\end{frame}

\begin{frame}{Today's Checklist}
\begin{enumerate}
  \item Identify players, strategies, and payoffs.
  \item Read the payoff matrix correctly.
  \item Use best responses to check Nash equilibrium.
  \item Remove dominated strategies when possible.
  \item For zero-sum games, try the saddle point test before mixed strategies.
\end{enumerate}
\end{frame}

\begin{frame}{Mini Quiz}
\begin{enumerate}
  \item In a Nash equilibrium, can the outcome be inefficient?
  \item In a zero-sum game, what does Player B want to do to Player A's payoff?
  \item What equality must hold for a saddle point?
  \item Why should we not compare the two payoffs inside the same cell?
\end{enumerate}
\end{frame}

\begin{frame}{Summary}
\begin{itemize}
  \item Game theory studies strategic decision making.
  \item Nash equilibrium means no player wants to deviate alone.
  \item Equilibrium can be stable but inefficient.
  \item Dominated strategies help simplify payoff tables.
  \item In zero-sum games, a saddle point gives a pure-strategy solution.
\end{itemize}
\vspace{0.3cm}
\begin{block}{Next}
Mixed strategies and graphical solution methods are natural next steps when no saddle point exists.
\end{block}
\end{frame}

\end{document}
